\begin{proof}[Proof of \cref{obj:thm:external-score-root-order}]
\emph{Step 1: Lower certificates.}
\Cref{obj:lem:external-two-source-lower-certificate} gives \(0<S_g\leq1\), the trial and score-log lower bounds, and the stated lower bound on the normalized risk. It also gives \(c_{\mathrm{log}}(g,\alpha)>0\). We construct an interval attaining the upper bound. % lean: CausalSmith.PartialID.UnlinkedPropensityAte.externalExcessRisk_lower_certificates

\emph{Step 2: Empirical distances and endpoint stability.}
Fix positive \(n,m\) and an admissible pair \((P,H)\). For each arm \(a\) and label \(r\), define the population finite measures
\[
\lambda_{ar}(B)=P(R=r,A=a,Y\in B),\quad B\subseteq[0,1],
\qquad
\sigma_{ar}(B)=\int_{B\cap C_r}p_a(e)\,H(de),\quad B\subseteq\mathcal E.
\]
Their total masses agree by the score-marginal, assignment, and release conditions. For finite measures \(\xi,\xi'\) on a compact interval \(\mathcal I\), write
\[
d_{\mathcal I}(\xi,\xi')
=
|\xi(\mathcal I)-\xi'(\mathcal I)|
+\int_{\mathcal I}|F_\xi(t)-F_{\xi'}(t)|\,dt,
\qquad F_\xi(t)=\xi((-\infty,t]).
\]
For a joint sample \(\omega\), set
\[
D_Y(\omega)=\sum_{a,r}d_{[0,1]}(\widehat\lambda_{ar}(\omega),\lambda_{ar}),
\qquad
D_E(\omega)=\sum_{a,r}d_{\mathcal E}(\widehat\sigma_{ar}(\omega),\sigma_{ar}).
\]
The empirical measures are those of \cref{obj:def:external-projection-handle}. Each empirical cumulative mass, including its total mass, is an average of independent variables in \([0,1]\). Its expected absolute deviation is at most \(1/(2\sqrt n)\) for a trial cell and \(1/(2\sqrt m)\) for a score-log cell. Integrating over intervals of lengths \(1\) and \(1-2\varepsilon\), respectively, and summing the \(2J\) cells gives
\[
\mathbb E_{\mathbb Q_{P,H}^{n,m}}D_Y\leq\frac{2J}{\sqrt n},
\qquad
\mathbb E_{\mathbb Q_{P,H}^{n,m}}D_E
\leq\frac{J(2-2\varepsilon)}{\sqrt m}. \label{eq:external-root-order-mean}
\]

Here is the stability estimate used below. For compatible finite measures \(\xi\) on \([0,1]\) and \(\zeta\) on \(\mathcal E\), let their common mass be \(\kappa\). When \(\kappa>0\), define
\[
\Psi_a^-(\xi,\zeta)
=\kappa\int_0^1 Q_{\xi/\kappa}(u)
 Q_{(p_a^{-1})_\#(\zeta/\kappa)}(1-u)\,du,
\qquad
\Psi_a^+(\xi,\zeta)
=\kappa\int_0^1 Q_{\xi/\kappa}(u)
 Q_{(p_a^{-1})_\#(\zeta/\kappa)}(u)\,du;
\]
when \(\kappa=0\), set both values to zero. For another compatible pair \((\xi',\zeta')\) of mass \(\kappa'\), put \(\kappa_\wedge=\min\{\kappa,\kappa'\}\). Suppose first that both masses are positive. The quantile transport identity on a compact interval \(\mathcal I\) identifies the integral of the absolute difference of normalized quantiles with the integral of the absolute difference of normalized CDFs. If \(\kappa\leq\kappa'\), then \(0\leq F_{\xi'}(t)\leq\kappa'\) gives, pointwise,
\[
\kappa\left|\frac{F_\xi(t)}{\kappa}-\frac{F_{\xi'}(t)}{\kappa'}\right|
\leq |F_\xi(t)-F_{\xi'}(t)|+|\kappa-\kappa'|.
\]
The same inequality with the pairs exchanged covers \(\kappa'\leq\kappa\). Applying it to the outcome and score CDFs and integrating over supports of lengths \(1\) and \(1-2\varepsilon\), respectively, yields
\[
\begin{aligned}
\kappa_\wedge\int_0^1
 |Q_{\xi/\kappa}(u)-Q_{\xi'/\kappa'}(u)|\,du
&\leq
\int_0^1|F_\xi(t)-F_{\xi'}(t)|\,dt+|\kappa-\kappa'|,\\
\kappa_\wedge\int_0^1
 |Q_{\zeta/\kappa}(u)-Q_{\zeta'/\kappa'}(u)|\,du
&\leq
\int_{\mathcal E}|F_\zeta(t)-F_{\zeta'}(t)|\,dt
 +(1-2\varepsilon)|\kappa-\kappa'|.
\end{aligned} \label{eq:external-root-order-transport}
\]
Since \(0\leq y\leq1\), \(p_a(e)\geq\varepsilon\), and \(e\mapsto p_a(e)^{-1}\) is \(\varepsilon^{-2}\)-Lipschitz, the product difference in either definition of \(\Psi_a^\pm\), multiplied by \(\kappa_\wedge\), is bounded using \cref{eq:external-root-order-transport} by
\[
\varepsilon^{-1}
 \left(\int_0^1|F_\xi-F_{\xi'}|+|\kappa-\kappa'|\right)
+\varepsilon^{-2}
 \left(\int_{\mathcal E}|F_\zeta-F_{\zeta'}|
       +|\kappa-\kappa'|\right). \label{eq:external-root-order-product}
\]
The remaining mass contributes at most \(\varepsilon^{-1}|\kappa-\kappa'|\). Because both distances contain that mass difference, \cref{eq:external-root-order-product} implies, with
\(\Lambda_\varepsilon=\varepsilon^{-1}+\varepsilon^{-2}\),
\[
|\Psi_a^\pm(\xi,\zeta)-\Psi_a^\pm(\xi',\zeta')|
\leq
\Lambda_\varepsilon
 \bigl(d_{[0,1]}(\xi,\xi')+d_{\mathcal E}(\zeta,\zeta')\bigr). \label{eq:external-root-order-stability}
\]
If either mass is zero, the bound \(0\leq\Psi_a^\pm(\xi,\zeta)\leq\varepsilon^{-1}\kappa\) and the mass term in the distance give \cref{eq:external-root-order-stability} directly. Summing \cref{eq:external-root-order-stability} over labels gives the same bound, with summed distances, for each arm endpoint. By \cref{obj:lem:fixed-released-law-baseline}, those endpoints at \((\lambda_{ar},\sigma_{ar})_{a,r}\) are \(\underline\mu_a,\overline\mu_a\). % lean: CausalSmith.PartialID.UnlinkedPropensityAte.externalJointTrial_meanDistance_le, CausalSmith.PartialID.UnlinkedPropensityAte.externalJointLog_meanDistance_le, CausalSmith.PartialID.UnlinkedPropensityAte.projectedCellSummand_abs_le, CausalSmith.PartialID.UnlinkedPropensityAte.projectedArmEndpoint_population_eq

\emph{Step 3: Honesty and excess length.}
Use the interval \(C_{n,m}\) of \cref{obj:def:external-projection-handle}. Its countable rational sieve and selected endpoints give measurable endpoints and a closed interval in \([-1,1]\). The release \(g:\mathcal E\to\mathcal R_J\), with nonempty \(\mathcal E\), also ensures \(J\geq1\). The radii in that construction are
\[
r_n=\frac{4J}{\alpha\sqrt n},
\qquad
r_m=\frac{2J(2-2\varepsilon)}{\alpha\sqrt m}.
\]
By \cref{eq:external-root-order-mean}, Markov's inequality applied to the nonnegative distances gives
\[
\mathbb Q_{P,H}^{n,m}\{D_Y<r_n,\ D_E<r_m\}\geq1-\alpha. \label{eq:external-root-order-coverage}
\]
On the strict-ball event in \cref{eq:external-root-order-coverage}, define the available slack
\[
\eta_* = \min\{r_n-D_Y(\omega),\,r_m-D_E(\omega)\}>0,
\qquad \eta_k=\frac{\eta_*}{k+1}\quad(k\geq1).
\]
For each arm-label cell, its population outcome and score measures have the same finite mass. Given \(\eta>0\), choose a common nonnegative rational mass and a common positive atom count for two equal-weight atomic approximations, with outcome and score atoms on their respective dense rational grids. To see the quantitative bound, first approximate each normalized measure by equally weighted quantile atoms; increasing the common atom count makes the sum of their integrated CDF errors arbitrarily small. Moving those atoms to the rational grids preserves that convergence. If the common population mass is \(\kappa\) and the chosen rational mass is \(\kappa'\), changing the weight of each atom from \(\kappa/N\) to \(\kappa'/N\) adds at most \((1+|\mathcal I|)|\kappa-\kappa'|\) to the distance on support \(\mathcal I\). Thus the combined weight error for the two supports is at most \((4-2\varepsilon)|\kappa-\kappa'|\). Choosing \(\kappa'\) sufficiently close to \(\kappa\) gives combined cell error below any prescribed tolerance; this also covers zero mass by taking rational mass zero.

Apply this construction in every cell with tolerance \(\eta_k/[2(J+1)]\). Since there are \(2J\) cells, it gives a compatible rational array \(b_k\) satisfying
\[
\sum_{a,r}\!\left(d_{[0,1]}(\lambda_{ar},\lambda_{ar}(b_k))
+d_{\mathcal E}(\sigma_{ar},\sigma_{ar}(b_k))\right)
<\frac{2J}{2(J+1)}\eta_k<\eta_k.
\]
Here \(\lambda_{ar}(b_k)\) and \(\sigma_{ar}(b_k)\) denote the two measures encoded by \(b_k\). Each of the two separate distance sums is therefore below \(\eta_k\). The triangle inequality gives empirical distances below \(D_Y(\omega)+\eta_k<r_n\) and \(D_E(\omega)+\eta_k<r_m\), so \(b_k\) belongs to both closed empirical balls. As \(k\to\infty\), \cref{eq:external-root-order-stability} makes all four arm endpoints converge to their population values. Write
\[
I(b_k)=[\underline\theta_k,\overline\theta_k],
\qquad
I(P_{\mathrm{rel}};H,g)=[\underline\theta_\star,\overline\theta_\star].
\]
Both candidate endpoints converge to their corresponding population endpoints. Fix any \(v\in I(P_{\mathrm{rel}};H,g)\) and clamp it to the candidate interval:
\[
v_k=\min\{\max\{v,\underline\theta_k\},\overline\theta_k\}\in I(b_k).
\]
If \(v<\underline\theta_k\), the distance \(|v_k-v|\) is at most \(|\underline\theta_k-\underline\theta_\star|\); if \(\overline\theta_k<v\), it is at most \(|\overline\theta_k-\overline\theta_\star|\); otherwise it is zero. Thus \(v_k\to v\). Every point of the population interval therefore lies in the closure of the union of candidate intervals, and hence in their closed convex hull. Write \(\theta(P)=\mathbb E_P[Y_1-Y_0]\). By \cref{obj:lem:fixed-released-law-baseline}, \(\theta(P)\) belongs to that sharp interval; its potential-outcome support also places it in \([-1,1]\). Thus the clipped hull contains \(\theta(P)\) on the event in \cref{eq:external-root-order-coverage}, proving
\[
C_{n,m}\in\mathcal J^{\mathrm{ext}}_{n,m,\alpha}(g),
\qquad
C_{n,m}(\omega)=\text{the closed interval of \cref{obj:def:external-projection-handle} at }\omega
\quad\text{for every }\omega. \label{eq:external-root-order-honesty}
\]

For any retained candidate, the two empirical-ball inequalities and the triangle inequality bound its summed distance from the population array by \(r_n+D_Y(\omega)+r_m+D_E(\omega)\). \Cref{eq:external-root-order-stability} therefore bounds each candidate arm endpoint's error by
\(\Lambda_\varepsilon(r_n+D_Y(\omega)+r_m+D_E(\omega))\). Each average-treatment-effect endpoint combines two arm endpoints; taking the hull and clipping it to \([-1,1]\) yields
\[
\left(\operatorname{length}(C_{n,m}(\omega))
-\operatorname{length}I(P_{\mathrm{rel}};H,g)\right)_+
\leq
4\Lambda_\varepsilon
 \bigl(r_n+D_Y(\omega)+r_m+D_E(\omega)\bigr). \label{eq:external-root-order-excess}
\]
In either fallback branch of the construction the interval is \(\{0\}\), so its positive excess is zero and \cref{eq:external-root-order-excess} holds there as well. Taking expectations and using \cref{eq:external-root-order-mean} gives, uniformly over admissible \((P,H)\),
\[
\mathbb E_{\mathbb Q_{P,H}^{n,m}}
 \left[\left(\operatorname{length}(C_{n,m})
-\operatorname{length}I(P_{\mathrm{rel}};H,g)\right)_+\right]
\leq
4\Lambda_\varepsilon
 \left(
 \frac{4J/\alpha+2J}{\sqrt n}
 +\frac{2J(2-2\varepsilon)/\alpha+J(2-2\varepsilon)}{\sqrt m}
 \right). \label{eq:external-root-order-mean-excess}
\]
The joint marginal and independence assumptions identify \(\mathbb Q_{P,H}^{n,m}\) with the product experiment in \cref{obj:def:external-experiment}. Hence honesty and \cref{eq:external-root-order-mean-excess}, followed by the infimum and supremum in \cref{obj:def:external-excess-risk}, give the same upper bound for \(\mathcal R^\star_{n,m,\alpha}(g)\). Set
\[
C_{\varepsilon,J,\alpha,g}
=
4\Lambda_\varepsilon
 \left(\frac{4J}{\alpha}+2J
       +\frac{2J(2-2\varepsilon)}{\alpha}
       +J(2-2\varepsilon)\right)+1>0.
\]
Both \cref{eq:external-root-order-mean-excess} and the minimax bound are then at most
\(C_{\varepsilon,J,\alpha,g}(n^{-1/2}+m^{-1/2})\), as required. % lean: CausalSmith.PartialID.UnlinkedPropensityAte.externalProjectionIntervalProcedure_honest, CausalSmith.PartialID.UnlinkedPropensityAte.externalProjectionIntervalProcedure_Icc, CausalSmith.PartialID.UnlinkedPropensityAte.externalProjectionIntervalProcedure_population_excessLength_le, CausalSmith.PartialID.UnlinkedPropensityAte.externalProjectionIntervalProcedure_population_meanExcessLength_le, CausalSmith.PartialID.UnlinkedPropensityAte.externalExcessRisk_le_of_honest_uniform_bound

\emph{Step 4: Normalized bounds.}
For positive \(n,m\), the upper bound just proved and \cref{obj:def:external-normalized-risk} imply
\[
0\leq T_{n,m,\alpha}(g)
=b_{n,m}\mathcal R^\star_{n,m,\alpha}(g)
\leq C_{\varepsilon,J,\alpha,g},
\qquad
b_{n,m}=(n^{-1/2}+m^{-1/2})^{-1}. \label{eq:external-root-order-normalized}
\]
The lower inequality follows from the positive trial certificate in \cref{obj:lem:external-two-source-lower-certificate}. Along the nonempty joint regime \(n,m\to\infty\), \(n/m\to\rho>0\), \cref{eq:external-root-order-normalized} bounds the limsup by \(C_{\varepsilon,J,\alpha,g}\) and ensures that the liminf is at most the limsup. The same cited lower certificate supplies
\[
\liminf T_{n,m,\alpha}(g)\geq
\max\left\{
\frac{c_{\mathrm{tr}}(\alpha)}{1+\sqrt\rho},
\frac{c_{\mathrm{log}}(g,\alpha)\sqrt\rho}{1+\sqrt\rho}
\right\}>0.
\]
Together with the score-log liminf bound in \cref{obj:lem:external-two-source-lower-certificate}, these are all the asserted inequalities. % lean: CausalSmith.PartialID.UnlinkedPropensityAte.externalExcessRisk_root_order
\end{proof}
